\section{Implementation order and falsifiable acceptance gates}
\label{sec:roadmap}

The right next step is not another solver family. It is to turn one of these
small replay computations into a checked source theorem. The B\"uchi rank lane
is the first choice: its checker is finite, its core proof uses natural-number
descent, and it exposes the important difference between typing a strategy and
proving its recurring behavior without initially requiring measure theory.

\begin{longtable}{@{}p{0.14\textwidth}p{0.34\textwidth}p{0.46\textwidth}@{}}
\toprule
Milestone & Deliverable & Acceptance gate \\
\midrule\endhead
1. Game replay & Lean arena definitions, both rank checkers, and both soundness theorems & Compile without placeholders; replay winning and losing fixtures; reject ownership, edge, start and rank mutations; record the actual theorem axiom inventory. \\
2. Source strategy & One explicit source arena bridge and a generated executable policy & Prove the original recurrent specification, not merely the finite table check. A source ownership change must invalidate the bridge or proof. \\
3. Pushdown replay & Productive summary and closed exclusion checkers, suffix theorem, and regular-target reduction & Prove an unbounded-stack safety example and a compressed positive reachability example without expanding their configuration spaces or traces. \\
4. Rational residuals & Stochastic-matrix validator and finite-horizon identities and bounds & Check the retry example and reject the harmonic-only self-loop. Verify costs and fractional terminal payoffs under their distinct meanings. \\
5. Probability bridge & A discrete source semantics linked to the matrix by exact pushforward & Prove almost-sure termination and the original expected-cost statement; changing an original probability must invalidate the proof. \\
6. Tactic integration & Contract recognition, search invocation, cached replay and source-goal closure & Evaluate against pinned baseline tactics and manually supplied model proofs, with model construction and kernel replay costs included. \\
\bottomrule
\end{longtable}

The comparison at milestone 6 should separate two sources of value. One track
supplies the same source bridge and local lemmas to all tactics, testing how
well the new worker finds the model certificate. A second track measures the
whole user experience, including contract recognition and bridge construction.
The first must not be described as an end-to-end automation result; the second
must not hide manually authored bridges as if the tactic synthesized them.
The existing repository's general benchmark discipline can be reused; only
these worker-specific distinctions need to be added.

\subsection{What should remain outside the first release}
General parity games, games with imperfect information, arbitrary temporal
formulas, stochastic games, unbounded-state Markov chains, and arbitrary
recursive Lean definitions should remain explicitly unsupported. The small
round-robin monitor, queue-based attractors, exact finite quotients, and
fraction-free linear solving are concrete extensions with clear proof
obligations. They are better next targets than a claim to solve all of the
larger subjects at once.

The three workers also cannot be combined merely by connecting their data
formats. An existential pop summary does not summarize an adversarial stack
game: its successful internal choices might belong to the opponent. A
probabilistic pushdown process can produce nonlinear probability equations and
is not automatically a finite rational Markov chain. Combining the present
workers requires a proved intermediate semantics, not a shared use of words
such as ``state,'' ``closure,'' or ``transition.''

\subsection{Assessment}
These extensions expand the kinds of statements Forge could automate rather
than only increasing a saturation budget. They give finite witnesses for
unbounded stack behavior, persistent strategies for infinite adversarial
execution, and properness-backed exact values for stochastic execution. Their
negative results have equally concrete meanings: a closed exclusion relation,
a counterstrategy bounding accepting visits, or a positive-probability path to
a nonterminal closed class.

The strongest completed deliverable is the combination of executable searches,
separate replay code, explicit source-bridge contracts, and exhaustive or
orthogonal tests on carefully delimited fragments. The principal unfinished
work is the Lean port and its end-to-end source examples. That work is now a
set of named checker theorems and translation obligations, rather than a
request to design a more powerful prover in the abstract.
